\documentclass[12pt,a4paper]{article}
\usepackage[utf8]{inputenc}
\usepackage[indonesian]{babel}
\usepackage{amsmath}
\usepackage{amsfonts}
\usepackage{amssymb}
\usepackage{geometry}

\geometry{margin=2.5cm}

\title{Menghitung Panjang Busur Kurva}
\author{Ahmad Fakhrul Bawani}
\date{}

\begin{document}

\maketitle

\section*{Soal}

Find the length of the following curve:
\begin{equation*}
  x = \frac{y^4}{4} + \frac{1}{8y^2}, \quad \text{dari } y = 1 \text{ sampai } y = 3
\end{equation*}

\subsection*{1. Rumus Umum}
Rumus untuk menghitung panjang busur (arc length) fungsi dalam variabel $y$ dari $y = c$ sampai $y = d$ adalah:
\begin{equation*}
  L = \int_{c}^{d} \sqrt{1 + \left(\frac{dx}{dy}\right)^2} \, dy
\end{equation*}

\subsection*{2. Menentukan Turunan Fungsi}
Mari tulis ulang fungsinya agar lebih mudah diturunkan:
\begin{equation*}
  x = \frac{1}{4}y^4 + \frac{1}{8}y^{-2}
\end{equation}

Cari turunan pertama $x$ terhadap $y$:
\begin{align*}
  \frac{dx}{dy} &= \frac{1}{4}(4y^3) + \frac{1}{8}(-2y^{-3}) \\
  \frac{dx}{dy} &= y^3 - \frac{1}{4y^3}
\end{align*}

Selanjutnya, kuadratkan hasil turunan tersebut:
\begin{align*}
  \left(\frac{dx}{dy}\right)^2 &= \left(y^3 - \frac{1}{4y^3}\right)^2 \\
  &= y^6 - 2(y^3)\left(\frac{1}{4y^3}\right) + \frac{1}{16y^6} \\
  &= y^6 - \frac{1}{2} + \frac{1}{16y^6}
\end{align*}

Tambahkan dengan 1 untuk dimasukkan ke dalam akar:
\begin{align*}
  1 + \left(\frac{dx}{dy}\right)^2 &= 1 + y^6 - \frac{1}{2} + \frac{1}{16y^6} \\
  &= y^6 + \frac{1}{2} + \frac{1}{16y^6}
\end{align*}

Perhatikan bahwa bentuk di atas merupakan kuadrat sempurna, yaitu:
\begin{equation*}
  y^6 + \frac{1}{2} + \frac{1}{16y^6} = \left(y^3 + \frac{1}{4y^3}\right)^2
\end{equation*}

\subsection*{3. Solusi Integrasi}
Substitusikan ke dalam rumus panjang busur dengan batas $c = 1$ dan $d = 3$:
\begin{align*}
  L &= \int_{1}^{3} \sqrt{\left(y^3 + \frac{1}{4y^3}\right)^2} \, dy \\
  L &= \int_{1}^{3} \left(y^3 + \frac{1}{4y^3}\right) \, dy \\
  L &= \int_{1}^{3} \left(y^3 + \frac{1}{4}y^{-3}\right) \, dy
\end{align*}

Lakukan integrasi:
\begin{align*}
  L &= \left[ \frac{1}{4}y^4 - \frac{1}{8y^2} \right]_{1}^{3} \\
  L &= \left( \left[ \frac{3^4}{4} - \frac{1}{8(3)^2} \right] - \left[ \frac{1^4}{4} - \frac{1}{8(1)^2} \right] \right) \\
  L &= \left( \left[ \frac{81}{4} - \frac{1}{72} \right] - \left[ \frac{1}{4} - \frac{1}{8} \right] \right)
\end{align*}

Samakan penyebut untuk mempermudah perhitungan:
\begin{align*}
  L &= \left( \frac{1458}{72} - \frac{1}{72} \right) - \left( \frac{2}{8} - \frac{1}{8} \right) \\
  L &= \frac{1457}{72} - \frac{1}{8}
\end{align*}

Samakan penyebut kedua pecahan menjadi 72:
\begin{align*}
  L &= \frac{1457}{72} - \frac{9}{72} \\
  L &= \frac{1448}{72}
\end{align*}

Sederhanakan dengan membagi pembilang dan penyebut dengan 8:
\begin{equation*}
  L = \frac{181}{9} = 20\frac{1}{9}
\end{equation*}

Jadi, panjang busur kurva tersebut adalah \textbf{$\dfrac{181}{9}$ atau $20\frac{1}{9}$ satuan panjang}.

\end{document}